% ============================================================================
%  Subdocumento 13. Introducción a las Innovaciones. Índice obligatorio del
%  Comunicado 10, secciones 8 y 11.
% ============================================================================

Este capítulo presenta la cartera de cinco innovaciones que exige el \art{28}{19}, una por cada tipo
y en el orden que fija el \comunicado{10}{sección 8}. Cada una responde a un problema medido del Caso.
El \art{30.2}{20} valora la pertinencia por sobre la novedad, y el \caso{capítulo 19}{44} pide
innovaciones pertinentes al transporte de carga y a los problemas de esta compañía, y no un catálogo de
tecnologías de moda. Ninguna de las cinco es una funcionalidad que las bases ya exigen, que es la
exclusión que declara el \art{30}{20}.

Las cinco comparten una estrategia: convierten en un dato medido algo que hoy se declara, se anota en
papel o se conoce tarde. La \tabref{13-cartera} resume el problema que resuelve cada una y el lugar
de la arquitectura del Subdocumento~4 donde se inserta.

\begin{tabla}{Cartera de innovaciones}{13-cartera}{R{8mm} L{44mm} X L{40mm}}
Tipo & Innovación & Problema del Caso que resuelve & Dónde se inserta \\ \midrule
1 & Portal del transportista con liquidación en curso & El 11 \% de las liquidaciones se corrige después de emitida & Portal y servicio de liquidación \\
2 & Despliegue sin detener la flota & La intervención a bordo solo ocurre cuando el camión pasa por un terminal & Actualización remota y terminales \\
3 & Semirremolque conectado & El sistema no sabe qué semirremolque va enganchado ni cuánto recorre & Equipo a bordo, portería y verificación bloqueante \\
4 & Esquema de adhesión y propiedad del dispositivo & Los dueños de camión no aceptan un equipo que delate cuándo trabajan para otros & Consentimiento, adhesión y equipo a bordo \\
5 & Bienestar del conductor y descanso en parador seguro & La alerta de jornada llega donde no hay dónde detenerse & Equipo a bordo y gobernanza de paradores \\
\end{tabla}
\notaTabla{Fuente: elaboración propia sobre el Caso, numerales 2.3, 4.4, 4.11 y 16.1, pp. 7, 10, 12 y 35, y capítulo 10, p. 23.}

Cuatro de las cinco usan el equipo a bordo que especifica el Formulario T-11, y dos usan el portal del
transportista. La \figref{13-carterafig} ubica cada innovación en las capas de la arquitectura.

\figura{figuras/13-innovaciones/cartera.png}{Cartera de innovaciones sobre las capas de la arquitectura, con los componentes que usa cada una}{13-carterafig}

En la figura, el equipo a bordo concentra las innovaciones 2, 3, 4 y 5, y el portal del transportista
las innovaciones 1 y 4. Esa concentración permite montar el valor nuevo sobre componentes que la
arquitectura ya necesita, sin agregar dispositivos al tractocamión.

Cada sección declara los siete elementos del \art{29}{19}: el problema, la tecnología, la madurez, el
diseño de la incorporación con su lugar en la arquitectura \transv*{RT-26.01}{44}, sus paquetes de la
estructura de descomposición del trabajo y su mes del cronograma \transv*{RT-26.02}{44}, el impacto
económico, el indicador y el riesgo de adopción. Los montos no aparecen en esta oferta técnica, porque
el \art{50.2}{29} lo prohíbe, y la valorización va en el Entregable 2 de la Oferta Económica, como
dispone el \art{30.1}{20}. Al menos una innovación es verificable en la marcha blanca, como valora
\transv{RT-26.08}{45}: la innovación 4 mide sus primeros anexos firmados desde el mes tres, y la 3 y
la 5 miden su indicador antes del mes 16.

El capítulo se apoya en el esquema de solución del Subdocumento~3, en la arquitectura del
Subdocumento~4, en el plan de trabajo del Subdocumento~7, que ubica los paquetes de cada innovación, y
en el plan de riesgos del Subdocumento~8. La innovación 1 usa además el motor de costeo del
Subdocumento~5. La 4 es contractual antes que técnica, y de su resultado depende cuántos de los 34
camiones de terceros sin equipo llegan a intervenirse. Las fichas completas van en el Formulario T-19, que acompaña
este capítulo como archivo propio (\texttt{AUDIT-Formulario-T-19.pdf}), conforme al
\comunicado{10}{sección 1}.

% ============================================================================
\section{Innovación 1 — Producto o servicio}\label{sec:13-innovacion-1}
% ============================================================================

La innovación 1 es de producto o servicio, según el \art{28}{19}. Su nombre es Portal del
transportista con liquidación en curso.

\paragraph{Problema u oportunidad} Hoy el transportista subcontratado espera nueve días a que ocho
personas calculen su liquidación, y el once por ciento de esos cálculos se corrige después
\caso*{numeral 4.11}{12}. Durante ese tiempo no tiene forma de saber cuánto va a cobrar ni por qué. La
innovación consiste en mostrarle su liquidación mientras se construye, viaje a viaje, en lugar de
entregársela cerrada al final del mes.

\paragraph{Tecnología que la sustenta} Portal segregado por identidad con vista propia de cada
transportista, alimentado por el mismo motor de costeo que produce la liquidación oficial. Cada viaje
aparece con su tarifa aplicada, sus tiempos de espera certificados por geocerca y sus descuentos, en el
momento en que el dato entra al sistema y no cuando el mes cierra.

\paragraph{Delimitación frente a las bases} Los criterios 21 y 29 obligan a que el transportista
autenticado consulte sus viajes, evidencias y liquidación en curso dentro del portal del mandante
\caso*{capítulo 18}{42}. Eso está comprometido en el alcance y no se presenta como innovación. Lo que
agrega esta ficha es que la evidencia salga del portal como un documento del transportista, verificable
por un tercero que no tiene acceso al sistema del mandante. Un expediente exportable con la jornada
acreditada de sus conductores, la vigencia de sus habilitaciones y la hoja de vida de sus equipos, que
él presenta a sus otros clientes, a la autoridad o a su aseguradora. Las bases no lo piden, y tampoco
cae en la exclusión del Capítulo 11 sobre administrar la contabilidad de los transportistas
\caso*{capítulo 11}{24}, porque no administra nada. Devuelve a su titular una evidencia ya producida.

\paragraph{Nivel de madurez} Escala de niveles de madurez tecnológica de uno a nueve
\parencite{iso16290}. La firma electrónica avanzada con verificación en línea está en nivel nueve, en
operación productiva y regulada en Chile desde 2002 \parencite{ley19799}, y es la línea base
comprometida. La credencial verificable alcanzó el estado de recomendación en mayo de 2025
\parencite{w3c_vcdm2} y se sitúa entre siete y ocho por su adopción todavía acotada en el ecosistema
logístico local.

\paragraph{Diseño de la incorporación} Un servicio de emisión de expediente en la capa de servicios de
negocio, que consume control de jornada, gestión documental y gestión de flota. Un servicio público de
verificación sin autenticación, que recibe un código y responde válido o inválido sin revelar
contenido. La firma la provee la bóveda de claves de la capa de seguridad. El portal expone la
solicitud y la descarga. La decisión entre firma avanzada y credencial verificable se documenta en el
mes cuatro con el levantamiento de destinatarios reales, la emisión se construye entre los meses diez y
doce, y el expediente queda disponible durante la marcha blanca de la Etapa~1. El motor de costeo del
que depende la liquidación en curso ya forma parte del alcance de la Etapa~1.

La construcción y despliegue del servicio se integran en los paquetes EDT 4.3 (Servicios de Liquidación y Trazabilidad) y EDT 7.3 (Portal del Transportista y Módulo de Expedientes), completando su desarrollo en el mes 12 para su verificación en la marcha blanca.

\paragraph{Impacto económico} La inversión es incremental sobre el portal ya presupuestado, agregando
el desarrollo de servicios y la suscripción de certificados de firma del emisor.

El efecto operacional reduce los costos administrativos de rectificación manual de liquidaciones y auditorías externas. En el flujo de caja de la Oferta Económica, la innovación se refleja en la partida de inversión CAPEX de desarrollo de software complementario y en los costos OPEX de infraestructura cloud para la custodia y verificación de credenciales.

\paragraph{Indicador de verificación} La \tabref{13-indicadores1} fija los indicadores. El primero
parte de cero porque la capacidad no existe, y el segundo usa la cifra del Caso.

\begin{tabla}{Indicadores de la innovación 1}{13-indicadores1}{X L{30mm} L{30mm} L{38mm}}
Indicador & Línea base & Meta & Momento de medición \\ \midrule
Transportistas adheridos que emiten al menos un expediente & 0 \% & 40 \% & Cierre de la marcha blanca, mes 15 \\
Liquidaciones corregidas después de emitidas & 11 \% & Menos de 2 \% & Tercer cierre mensual tras el paso a producción \\
\end{tabla}
\notaTabla{Fuente: elaboración propia sobre el Caso, numeral 4.11, p. 12.}

Las dos metas se miden con datos del propio portal y del motor de costeo. Se espera además una
reducción del tiempo de cierre mensual.

\paragraph{Riesgo de adopción} Que el transportista no perciba utilidad y no emita el expediente, con
probabilidad media e impacto alto. Se mitiga asistiendo la primera emisión en el terminal durante el
enrolamiento. El riesgo secundario es que los destinatarios no acepten el documento como prueba,
mitigado levantando destinatarios reales en el mes cuatro antes de construir.

En caso de baja adopción o retraso en la emisión por parte de los transportistas, el plan de contingencia contempla la generación automática del expediente consolidado al cierre de cada ciclo mensual y su envío asistido por correo electrónico certificado, garantizando la disponibilidad de la evidencia sin alterar los flujos de liquidación estándar.

% ============================================================================
\section{Innovación 2 — Proceso}\label{sec:13-innovacion-2}
% ============================================================================

La innovación 2 es de proceso, según el \art{28}{19}. Su nombre es Despliegue continuo y modular sin detención de flota.

\paragraph{Problema u oportunidad} Un camión detenido no produce \caso*{capítulo 10, restricción 10}{24},
la intervención física en cabina solo puede realizarse cuando la unidad ingresa a un terminal, y ese
paso ocurre en promedio cada seis días mientras que el 22 por ciento de la flota subcontratada ingresa
menos de una vez al mes \caso*{numeral 2.3}{7}. La innovación transforma el despliegue de hardware: deja
de ser un proyecto de detención masiva que paraliza la operación comercial y pasa a ser un proceso
continuo, no intrusivo, que aprovecha las ventanas naturales de estadía en terminal mediante kits
modulares preconfigurados de sustitución rápida.

\paragraph{Tecnología que la sustenta} Gestión centralizada y aprovisionamiento remoto del parque
telemático mediante Azure IoT Hub Device Update y Device Twins, configuración precargada en fábrica y
kits de reemplazo con arneses estandarizados de acople rápido por familia de tractocamión (Scania,
Volvo, Mercedes-Benz, Freightliner). Anticipación de llegada a terminal a partir de la posición de la
propia flota, de modo que el equipo técnico en terreno disponga del kit específico en el andén antes de
que el camión ingrese a patio.

\paragraph{Delimitación frente a las bases} El alcance comprende el despliegue directo de los
dispositivos provistos por el mandante en las 148 unidades propias y en las 34 unidades sin equipo
cuyos dueños suscriban el convenio de comodato de la innovación 4. No interviene físicamente los 192
camiones de terceros que ya cuentan con dispositivos telemáticos preexistentes, los cuales se integran
de manera lógica mediante adaptadores de interoperabilidad sin alterar su hardware privado, respetando la restricción 3 \caso*{capítulo 10}{23}. La
delimitación frente a las bases radica en que el Caso exige disponer de stock de reemplazo, pero no
resuelve el método operativo para instalar sin detener el flete comercial ni define la logística de
intervención en tránsito.

\paragraph{Diseño de la incorporación} La intervención se ejecuta en los cinco terminales operativos
(San Bernardo, Antofagasta, Talca, Los Ángeles y Puerto Montt) durante las ventanas de relevo de
conductor o pausas de carga y descarga, con un procedimiento estandarizado cronometrado en 30 minutos o
menos por vehículo. Tras el montaje, el dispositivo ejecuta una rutina automatizada de autodiagnóstico
y enrolamiento (autenticación mutua por certificado, enlace celular, fijación GNSS y lectura de bus CAN) que valida la
correcta operación en la barrera de salida sin requerir instrumentación adicional en patio. La
actividad se articula en el paquete EDT 4.2 (Homologación e Instalación de Búfer a Bordo, meses 4 a 8)
y el EDT 4.5 (Integración Telemática y Lectores CAN sin contacto, meses 6 a 9). El levantamiento de las
familias de conectores y la distribución de cabina de los tractocamiones propios se consolida en el
piloto cronometrado de los primeros 10 camiones durante el mes 4.

Estos paquetes corresponden a EDT 4.2 (Instalación de Dispositivos a Bordo y Búfer) y EDT 4.5 (Integración Telemática y Lectores CAN sin contacto), estructurados para ejecutarse en ventanas de patio entre los meses 4 y 9.

\paragraph{Nivel de madurez} Escala de madurez tecnológica ISO 16290:2013 \parencite{iso16290}. Los procedimientos de reemplazo modular de hardware y aprovisionamiento telemático mediante Device Twins alcanzan un nivel TRL 8, ampliamente probado en operaciones logísticas de alta exigencia.

\paragraph{Impacto económico} La inversión requerida se concentra en los arneses de acople rápido y el stock de búferes de intercambio. En costo operacional genera ahorros inmediatos al suprimir las horas de lucro cesante por inmovilización de camiones. Se refleja en la partida de inversión de equipamiento y logística de terreno de la Oferta Económica.

\paragraph{Indicador de verificación} Cero horas de lucro cesante o inmovilización de flota
comercialmente activa atribuible a la instalación. Ritmo sostenido de intervención de 35 a 45
tractocamiones mensuales distribuidos en la red de terminales, alcanzando el 100 \% de la flota propia
(148 unidades) al cierre del mes 8. Reducción del tiempo medio de sustitución en terreno a menos de
30 minutos por camión.

\paragraph{Riesgo de adopción} Riesgo de dispersión geográfica imprevista de unidades subcontratadas que retrase su paso por terminal. Se mitiga mediante el seguimiento telemático de proximidad y la programación de turnos móviles de instalación en talleres regionales. La contingencia contempla la intervención coordinada durante detenciones obligatorias de mantención preventiva en talleres autorizados.

% ============================================================================
\section{Innovación 3 — Tecnológica o de arquitectura}\label{sec:13-innovacion-3}
% ============================================================================

La innovación 3 es de tipo tecnológico o de arquitectura, según el \art{28}{19}. Su nombre es
Semirremolque conectado: una baliza Bluetooth en cada semirremolque propio permite que el sistema sepa
qué semirremolque va realmente enganchado, cuántos kilómetros recorre y en qué condición viaja la carga
refrigerada.

\paragraph{Problema u oportunidad} La compañía tiene 210 semirremolques propios de cinco tipos: rampla
plana, furgón seco, furgón refrigerado, tolva y portacontenedores \caso*{numeral 2.2}{6}. En tres años
llegan a 245 \caso*{numeral 14.1}{29}. De ellos, 44 son equipos refrigerados que concentran su
actividad entre diciembre y abril \caso*{numeral 2.2}{6}. La revisión técnica rige para los 210
semirremolques \caso*{capítulo 12}{25} y es una de las vigencias que la verificación bloqueante
comprueba antes de asignar un viaje \caso*{numeral 4.4}{10}. Las bases piden esa comprobación, pero no
dicen cómo sabe el sistema qué semirremolque va enganchado, y con una declaración del despachador el
requisito se cumple. Si la declaración es errónea, la verificación aprueba la vigencia de un
semirremolque que no es el que sale.

El mantenimiento tiene el mismo vacío. El taller atiende 148 tractocamiones y 210 semirremolques con
un plan preventivo por kilometraje, y el kilometraje se lee del odómetro cuando el camión pasa por el
taller \caso*{capítulo 8}{19}. Cuando un semirremolque cambia de tractocamión, el odómetro del camión
deja de reflejar su recorrido. El criterio 26 pide que el mantenimiento preventivo se gatille con
kilometraje real y no con una estimación \caso*{capítulo 18}{43}.

\paragraph{Tecnología que la sustenta} Balizas Bluetooth de baja energía Teltonika EYE Sensor BTSMP1,
con identificador, temperatura, humedad, acelerómetro y sensor magnético de puerta, protección IP67,
operación de $-20$ a $+60$~°C y batería de cinco años o más con una emisión cada 10 segundos
\parencite{teltonika_eye}. La baliza se fija en el chasis delantero, a la sombra y sin cables. La lee
el iWave G26I, que ya trae Bluetooth 5.0 \parencite{iwave_g26i}, y un lector Minew G1 en la portería de
cada terminal \parencite{minew_g1}. El equipo a bordo infiere el acople cuando la señal de la baliza es
fuerte y estable y su acelerómetro se mueve junto con el camión, práctica que el fabricante describe
para el seguimiento de remolques \parencite{teltonika_trailer}. El lector de portería registra el
semirremolque que sale aunque lo tire un camión sin equipo audIT.

\paragraph{Delimitación frente a las bases} El seguimiento de remolques con balizas ya se vende en el
mercado, así que la innovación no está en la baliza. Está en tres usos que cambian decisiones del caso.
La verificación bloqueante comprueba el semirremolque que efectivamente sale y no el declarado. El
kilometraje de cada semirremolque sale de su tiempo acoplado sobre la traza del tractocamión, lo que
extiende a los semirremolques el kilometraje real que el criterio 26 pide y que la solución base mide
solo en el camión. Y los 44 equipos refrigerados registran temperatura y apertura de puerta durante el
viaje. Ninguno de los catorce asuntos que la arquitectura debe resolver \caso*{numeral 17.4}{39} ni
ninguna de las decisiones pendientes \caso*{numeral 16.1}{34} pide identificar el semirremolque
acoplado ni medir la condición de la carga, de modo que la innovación no cae en la exclusión del
\art{30}{20}.

\paragraph{Nivel de madurez} Escala de niveles de madurez tecnológica de uno a nueve
\parencite{iso16290}. Las balizas y su lectura por Bluetooth están en nivel 9, porque son productos
comerciales en operación \parencite{teltonika_eye,minew_g1}. La inferencia del acople y su uso dentro
de la verificación bloqueante parten en nivel 6 y llegan a nivel 7 cuando se demuestran en la operación
durante la marcha blanca, como pide \transv{RT-26.03}{44}.

\paragraph{Diseño de la incorporación} La \figref{13-semirremolquefig} muestra dónde se inserta la
innovación en la arquitectura.

\figuraHorizontal{figuras/13-innovaciones/semirremolque.png}{Semirremolque conectado: la baliza, sus dos lectores y los tres servicios de negocio que usan el acople}{13-semirremolquefig}

La baliza vive en el semirremolque y no agrega nada al tractocamión. En la capa de borde, el G26I y el
lector de portería la leen por Bluetooth y envían el evento a IoT Hub, el primero por la red celular y
el segundo por la red del terminal. En la capa de integración, Event Hubs distribuye el evento a tres
servicios. El contexto Flota y activos usa el identificador del semirremolque acoplado en la
verificación bloqueante y suma sus kilómetros para el plan preventivo. El contexto Telemetría y
geocercas guarda la temperatura y la apertura de puerta en la base de series de tiempo. La innovación
consume el flujo de telemetría de IoT Hub y expone a la verificación bloqueante el identificador del
semirremolque acoplado, sin interfaces nuevas hacia terceros (secciones 4.1 y 4.2). El Formulario T-11
declara el equipo: 210 balizas más 21 de reposición, 35 más por el crecimiento a 245 semirremolques y
cinco lectores de portería.

Las balizas se instalan en la Etapa~1 cuando cada semirremolque pasa por un terminal. Los 44
refrigerados se intervienen fuera de la temporada de diciembre a abril, en que la operación no admite
intervenciones \caso*{numeral 13.2}{27}. La verificación con el semirremolque medido se activa en la
marcha blanca, y su primera medición ocurre antes del mes 16.

La ejecución se descompone en los paquetes EDT 3.4 (Instalación de Balizas Bluetooth en Semirremolques y Lectores de Portería, meses 4 a 9) y EDT 4.4 (Integración con Verificación Bloqueante y Despacho, meses 8 a 11), completando la certificación de acople previo al hito de marcha blanca.

\paragraph{Impacto económico} La inversión es el equipamiento que compra el mandante, conforme al
\caso{capítulo 11}{24}: las balizas y los lectores de portería del Formulario T-11, más la
integración en los tres servicios de negocio. En costo operacional agrega la reposición de balizas al
término de su batería y un tráfico celular menor, porque el evento de acople viaja en los mensajes que
el G26I ya envía. El beneficio tiene tres partes: impedir la salida de un semirremolque con la
revisión técnica vencida, programar la mantención de los semirremolques con kilometraje medido en lugar
de estimado, y respaldar ante el cliente la condición de la carga refrigerada. El flujo de caja de la
Oferta Económica refleja la innovación como inversión en equipamiento en la Etapa~1 y como costo de
reposición durante la operación.

\paragraph{Indicador de verificación} La \tabref{13-indicadores3} fija los indicadores. Las líneas
base describen lo que el sistema actual no registra, sin cifras que el Caso no entregue.

\begin{tabla}{Indicadores de la innovación 3}{13-indicadores3}{X L{32mm} L{30mm} L{36mm}}
Indicador & Línea base & Meta & Momento de medición \\ \midrule
Asignaciones con semirremolque propio verificado por medición & No se registra & 95 \% o más & Marcha blanca, antes del mes 16 \\
Semirremolques propios con kilometraje medido & No se registra & 100 \% de los que tienen baliza & Tercer mes de operación \\
Viajes refrigerados con registro continuo de temperatura & No se registra en el sistema & 100 \% & Primera temporada de diciembre a abril \\
\end{tabla}
\notaTabla{Fuente: elaboración propia sobre el Caso, numerales 2.2 y 4.4, pp. 6 y 10, y capítulo 18, p. 43.}

La primera meta no llega al 100 \% porque un semirremolque propio tirado por un camión sin equipo
audIT solo se lee en la portería. Los tres indicadores se miden con datos que la propia solución
produce, de modo que su verificación no depende de terceros.

\paragraph{Riesgo de adopción} La \tabref{13-riesgos3} declara los riesgos con su probabilidad, su
impacto, la mitigación y la contingencia, como pide \transv{RT-26.04}{44}.

\begin{tabla}[espacio=22]{Riesgos de adopción de la innovación 3}{13-riesgos3}{L{30mm} L{16mm} L{14mm} X X}
Riesgo & Probabilidad & Impacto & Mitigación & Contingencia \\ \midrule
Acople falso con la baliza de un semirremolque vecino en el patio & Media & Alto & Exigir señal estable y movimiento conjunto en los primeros minutos de marcha & La verificación vuelve a la declaración con la lectura de portería como control \\
Temperatura bajo $-20$~°C en el paso cordillerano & Baja & Medio & Montaje a la sombra del chasis y prueba en el invierno de la marcha blanca & Un modelo de mayor rango térmico en los semirremolques que cruzan \\
Suplantación del identificador de la baliza & Baja & Alto & El acople exige movimiento conjunto y la lectura de portería & El acople se trata como dato complementario y no como única prueba \\
Batería que dura menos de lo declarado & Baja & Bajo & Emisión cada 10 segundos y vigilancia del nivel que reporta la baliza & Reposición desde la reserva del 10 \% del Formulario T-11 \\
\end{tabla}
\notaTabla{Fuente: elaboración propia sobre \textcite{teltonika_eye}.}

El riesgo mayor es el acople falso, porque afecta a la verificación bloqueante, y por eso su
contingencia devuelve el control a la declaración sin detener la operación.

En cumplimiento de RT-26.07, la incorporación de balizas BLE fue evaluada en el modelado de amenazas STRIDE de la sección 4.1, incorporando mitigaciones criptográficas con tokens rotatorios de autenticación para neutralizar riesgos de suplantación (\textit{spoofing}) o manipulación en patio.

% ============================================================================
\section{Innovación 4 — Modelo de negocio o de contratación}\label{sec:13-innovacion-4}
% ============================================================================

La innovación 4 es de modelo de negocio o de contratación, según el \art{28}{19}. Su nombre es
Esquema de adhesión y propiedad del dispositivo, y responde a lo que el \caso{capítulo 19}{44} pide
para este tipo: hacerse cargo de la relación con los transportistas subcontratados.

\paragraph{Problema u oportunidad} Financiamiento del dispositivo, incentivos de adhesión,
consentimiento granular y beneficios verificables. La primera pregunta que hace un dueño de camión es
de quién es el equipo, quién lo paga y qué pasa con él si deja de trabajar con la compañía. La segunda
es si el aparato va a delatar cuándo trabaja para la competencia. Mientras esas dos preguntas no tengan
respuesta, ningún diseño técnico se despliega. La innovación separa la propiedad del equipo de la
propiedad del dato, y le entrega la segunda al dueño del camión.

\paragraph{Tecnología que la sustenta} Figura de comodato sobre el equipamiento, con condiciones de
retiro y de traspaso definidas desde el inicio. Modo de privacidad implementado en el firmware del
dispositivo y no en configuración de servidor, de modo que la posición deja de emitirse hacia el
mandante cuando el camión trabaja para otro cliente, y esa garantía sea verificable por el dueño en
lugar de prometida. Consentimiento granular y revocable administrado desde el portal, con registro
auditable de cada acceso a la información de localización.

\paragraph{Delimitación frente a las bases} El Capítulo 11 resuelve quién compra el hardware
\caso*{capítulo 11}{24}, y presentar esa regla como innovación sería presentar como propia una decisión
que ya está en las bases. Lo que el Caso deja abierto es la decisión quinta del numeral 16.1, sobre el
dispositivo instalado en un camión de un tercero \caso*{numeral 16.1}{34}, y la restricción 2, que
obliga a conseguir por contrato, por incentivo o por diseño lo que dependa de terceros
\caso*{capítulo 10}{23}. Esta ficha agrega tres cosas que las bases no piden. La ventana de transmisión
gobernada en el firmware, que convierte una promesa de privacidad en una propiedad verificable. El
comodato con retiro sin costo, que elimina el riesgo patrimonial del transportista sobre un activo que
es suyo. Y el financiamiento del incentivo con el recupero de un ingreso que hoy se pierde, de modo que
la adhesión no compita con el margen operacional de nueve por ciento del mandante
\caso*{numeral 2.3}{7}.

\paragraph{Nivel de madurez} Los componentes técnicos se evalúan en la escala de niveles de madurez
tecnológica \parencite{iso16290}. El control de transmisión por ventana en el dispositivo está en nivel
ocho, porque la gestión remota de configuración de equipos conectados es tecnología productiva y lo
específico es la regla de negocio que la gobierna. El registro de consentimiento granular y revocable
está en nivel ocho y su exigencia es normativa \parencite{ley21719}. El componente contractual no se
califica en esta escala porque no es una tecnología. El comodato de equipamiento a proveedores de
servicio es figura de uso corriente y no requiere desarrollo. El reparto de recupero sobre evidencia
aportada tiene precedente acotado y poca documentación pública en el sector nacional.

\paragraph{Diseño de la incorporación} La ventana se aplica en la unidad telemática y en el búfer
local, que almacenan fuera de ella y no transmiten. El concentrador de dispositivos distribuye la
configuración de ventana a cada equipo. Un servicio de consentimiento y un servicio de adhesión, ambos
nuevos, definen la ventana desde la asignación del viaje y registran la revocación. El portal expone la
consola de permisos con bitácora. La adhesión comienza en el mes uno, antes que cualquier construcción,
porque el numeral 13.1 advierte que hay decisiones cuyo plazo no lo fija la tecnología sino una
negociación con terceros, y las negociaciones no se paralelizan \caso*{numeral 13.1}{26}. La cohorte
piloto va entre los meses seis y nueve, la consola entre el nueve y el doce, y el resultado se mide en
la marcha blanca. El tratamiento contable del comodato y el reparto concreto del recupero de
sobreestadía requieren definición conjunta con el mandante.

Las actividades se articulan en los paquetes EDT 1.3 (Validación Jurídica del Convenio de Comodato y Privacidad, meses 1 a 3), EDT 5.2 (Campaña de Enrolamiento y Adhesión en Terminales, meses 3 a 8) y EDT 7.4 (Consola de Consentimiento y Soberanía del Dato en Portal, meses 6 a 10).

\paragraph{Impacto económico} La adquisición del equipamiento es del mandante conforme al
Capítulo 11, y audIT especifica la cantidad sobre las poblaciones del Caso, con los 34 camiones de
terceros sin equipo incluidos (sección 4.2.1). Las partidas propias son el diseño y la validación
jurídica del anexo, la campaña de enrolamiento en cinco terminales con presencia en horario de relevo,
y la consola de consentimiento, todas valorizadas en el flujo de caja de la Oferta Económica. En costo
operacional aumentan la conectividad, el soporte y la reposición del parque en comodato, y disminuyen
el costo de la liquidación mensual y el de gestionar las objeciones de cobro. El beneficio se apoya en
una base verificada: los clientes objetaron el setenta y uno por ciento de lo facturado en 2025 por
tiempos de espera \caso*{numeral 4.7}{11}, porque la hora de llegada se anota en papel. La meta de
reducir la objeción bajo el veinte por ciento genera un recupero que se valoriza en la Oferta Económica
y no se compromete como ingreso. Ese recupero es el que financia el incentivo, y su reparto concreto se
fija en el anexo contractual.

En la estructura de costos de la Oferta Económica, la inversión inicial cubre la asesoría legal y la campaña de difusión en terreno, mientras que el flujo de caja operacional refleja la amortización del comodato y el retorno generado por la disminución drástica de objeciones comerciales en tiempos de espera.

\paragraph{Indicador de verificación} Línea base cero de ciento cuarenta y ocho transportistas
adheridos. Meta de setenta por ciento al cierre de la Etapa~1 y noventa por ciento al cierre de la
Etapa~2. La primera medición de anexos firmados ocurre desde el mes tres, mucho antes del paso a
producción, lo que satisface con holgura el requisito deseable de que al menos una innovación sea
verificable durante la marcha blanca. Indicadores complementarios: tasa de revocación del
consentimiento bajo el diez por ciento de los adheridos, y objeción sobre cobros de espera respaldados
bajo el veinte por ciento desde el setenta y uno actual.

\paragraph{Riesgo de adopción} Que la adhesión no alcance el setenta por ciento en la Etapa~1, con
probabilidad media e impacto alto, porque limita jornada, posición y emisiones sobre el sesenta coma
cuatro por ciento de la capacidad. Se mitiga comenzando en el mes uno y mostrando beneficio verificable
antes de pedir el equipo, y la contingencia es escalonar el incentivo y extender la modalidad de datos,
que no requiere instalar nada. Que el transportista desconfíe de la ventana de transmisión, que se
mitiga haciéndola auditable por él mismo desde su consola de permisos. Que los clientes no acepten la
evidencia telemática como respaldo del cobro, lo que eliminaría la fuente de financiamiento del
incentivo, con contingencia de financiarlo con la reducción del costo de liquidación, que no depende
del cliente. Y que el anexo no resista revisión jurídica, con probabilidad baja e impacto alto,
mitigado por la validación entre los meses uno y tres, antes de construir.

% ============================================================================
\section{Innovación 5 — Experiencia de usuario, sostenibilidad o impacto social}\label{sec:13-innovacion-5}
% ============================================================================

La innovación 5 es de experiencia de usuario, sostenibilidad o impacto social, según el \art{28}{19}.
Su nombre es Bienestar del conductor y descanso en parador seguro.

\paragraph{Problema u oportunidad} Una alerta que avisa al conductor cuando su cuota de cinco horas
continuas está por agotarse \parencite{codigo_trabajo_25bis}, entregada en medio de una cuesta sin
bermas o en un tramo desértico sin paraderos habilitados para combinaciones de 18,6 metros y 45
toneladas, no soluciona nada: lo fuerza a detenerse en bermas peligrosas o a seguir rodando en
infracción. La innovación abandona el temporizador pasivo e implementa un sistema sociotécnico de
acompañamiento predictivo del descanso, que cruza la cinemática del camión con la curva biológica de
alerta circadiana y un catálogo georreferenciado y calificado de paradores de carga pesada.

\paragraph{Tecnología que la sustenta} Motor predictivo embebido en el equipo a bordo de la
sección~4.2.2, que evalúa en tiempo real el tiempo de conducción acumulado, la
topografía del tramo siguiente y la ventana de mínima alerta circadiana (\emph{Window of Circadian
Low}, de 02:00 a 06:00~h). Catálogo georreferenciado local en SQLite, en modo de registro anticipado de escritura, sobre la
memoria flash industrial de 8~GB del equipo, especificada para las condiciones reales de uso que pide
\transv{RT-08.11}{19}. El catálogo califica
capacidad geométrica de estacionamiento, servicios de higiene, duchas, agua potable y seguridad
perimetral. Enclavamiento cinético estricto de pantalla ante cualquier movimiento del vehículo, con
velocidad sobre 0~km/h o freno de mano liberado, conforme a la Ley N.º~21.377 \parencite{ley21377}.
Las sugerencias llegan por síntesis de voz local en español por el altavoz del vehículo, que es el
canal de la alerta de jornada \caso*{RT-16.21}{33}, sin exigir desviar la mirada ni admitir
interacción táctil en marcha \caso*{RT-13.08}{33}. En modo detenido, interfaz ergonómica para guantes gruesos, con botones
de 60~×~60~mm o más e iluminación ámbar de alto contraste.

\paragraph{Delimitación frente a las bases} El cálculo reactivo de distancia al lugar seguro de
detención \caso*{RT-09.01}{32}, la autenticación sin manipulación \caso*{RT-12.11}{32}, la interfaz
para guantes sin interacción en marcha \caso*{RT-13.08}{33} y la resiliencia desconectada de 72 horas
\caso*{RT-17.01}{33} constituyen el alcance base comprometido y no se presentan como
innovación. Lo que agrega esta ficha son tres elementos. La anticipación predictiva circadiana, que
sugiere pausas preventivas antes de que el conductor ingrese a tramos críticos sin infraestructura o
enfrente el valle biológico de microsueño de madrugada. En esa franja, a las 04:40, ocurrió el
accidente del 14 de febrero de 2026, con un conductor que no había completado su descanso
\caso*{capítulo 1}{4}. La ontología, gobernanza y sincronización por actualización remota en terminales de un catálogo calificado para
camiones de 45 toneladas con servicios dignos de descanso. Y un modelo social y ergonómico no punitivo
para los 258 conductores externos subcontratados \caso*{numeral 2.3}{6}, conforme a las restricciones 1
y 2 \caso*{capítulo 10}{23} y a la Ley N.º~20.123 \parencite{ley20123}, que transforma una herramienta de control patronal en
un asistente de seguridad y bienestar.

\paragraph{Nivel de madurez} Escala de madurez tecnológica de uno a nueve \parencite{iso16290}. El
hardware telemático de borde, el almacenamiento local, la síntesis de voz local y el enclavamiento
cinético están en nivel ocho, con componentes comerciales de uso maduro en flotas pesadas. El modelo
predictivo circadiano y la ontología de paradores adaptada a la red vial chilena se sitúan en nivel
siete, demostrados en entorno operacional representativo de la Ruta 5 y pasos cordilleranos, apoyados
en estándares de ergonomía vehicular \parencite{iso15005,iso9241_210}, en investigación sobre fatiga en
el transporte de carga \parencite{nasem2016_fatiga} y en la normativa vial y laboral nacional
\parencite{ley21377,codigo_trabajo_25bis}.

\paragraph{Diseño de la incorporación} Se articula en la capa de borde (motor predictivo a bordo,
catálogo local y síntesis de voz), en la capa de servicios (servicio de gobernanza de paradores y
seguridad vial en la torre de control) y en la de presentación (portal para consulta pasiva). Su
trazabilidad comprende los paquetes EDT 2.4 (Diseño Ergonómico de la Interfaz y Protocolo de Guantes,
meses 4 a 7), EDT 3.7 (Catálogo de Paradores en la Campaña de Cobertura, meses 3 a 6),
EDT 4.6 (Software de Enclavamiento Cinético y Alerta Circadiana, meses 7 a 10) y EDT 7.2 (Validación
Operacional en Marcha Blanca con Conductores Reales, meses 13 a 15). El beneficio es verificable antes
del mes 16, como valora \transv{RT-26.08}{45}. La validación de los paradores y sus atributos se
consolida durante la campaña de cobertura de la Etapa~1, y la calibración de la ventana circadiana con
los turnos reales se ajusta en los talleres de diseño ergonómico del mes cuatro.

La numeración y alcance de los paquetes EDT 2.4, 3.7, 4.6 y 7.2 se encuentran plenamente alineados con la estructura de descomposición del trabajo del Subdocumento 7.

\paragraph{Impacto económico} La inversión requerida es incremental sobre el software de borde y los
talleres ergonómicos participativos, mientras que el levantamiento del catálogo se absorbe dentro del
recorrido de la campaña de medición de cobertura que exige el \caso{RT-03.24}{31}, sin costo
logístico adicional. En costo operacional
reduce la búsqueda errática de estacionamiento en ruta, con un efecto en combustible que se mide en la
Etapa~1. El beneficio central radica en la mitigación del riesgo de siniestros graves por microsueño y
en la supresión de multas por la Ley N.º~21.377 y el artículo 25 bis.

En el flujo de caja de la Oferta Económica, la innovación se registra en la partida CAPEX de desarrollo de software de borde y talleres ergonómicos, y en OPEX bajo la línea de servicios de actualización y soporte de cartografía segura.

\paragraph{Indicador de verificación} Línea base cero de alertas con parador calificado garantizado,
con meta de 98 \% o más de efectividad en ruta nacional, medida mensualmente en la marcha blanca. Cero
eventos de interacción táctil con el vehículo en movimiento, sobre 0~km/h, con meta de 100 \% de
bloqueo cinético.
Carga cognitiva NASA-TLX con guantes bajo 35 puntos, con línea base que se mide en los talleres de la
Etapa~1. Y adherencia voluntaria a paradas seguras de los conductores externos de 85 \% o más, con
línea base que se mide en la misma etapa.

\paragraph{Riesgo de adopción} Que los 258 conductores externos perciban el sistema como vigilancia
patronal encubierta, con probabilidad media a alta e impacto alto. Se mitiga mediante diseño
ergonómico participativo en horario de relevo en los cinco terminales y orientando la herramienta al
bienestar, con recomendación de sitios con duchas, seguridad y convenios comerciales. La contingencia
es que la unidad a bordo conmute a notificación puramente acústica por el altavoz del vehículo, sin
requerir teléfonos personales ni botones adicionales en cabina, manteniendo la seguridad activa sin
fricción sindical.
